\documentclass[10pt,a4paper]{amsart}
\usepackage{amsmath,amssymb,amsthm,graphicx}
% Retained preamble for the Stacks Project chapter divisors.tex (Divisors).
% The chapter file pinned-original-source.tex carries only \input{preamble} on its first line, so no
% preamble text lives inside the chapter itself.  The Stacks Project uses one single shared preamble for
% every chapter; the snapshot stages it verbatim in sources/stacks-chow/preamble.tex.  The block below is
% copied mechanically, line for line, from that file's lines 14-18 (the Xy-pic setup) and lines 154-217
% (the native theorem declarations together with the mathematical macros).  Nothing was added, removed or
% redefined.  The wrapper class and the amsmath/amssymb/amsthm/graphicx base packages are supplied by the
% assembler, so the class/\externaldocument part of the shared preamble is not repeated here.
% For commutative diagrams we use Xy-pic
\usepackage[all]{xy}

% We use 2cell for 2-commutative diagrams.
\xyoption{2cell}
% Theorem environments.
%
\theoremstyle{plain}
\newtheorem{theorem}[subsection]{Theorem}
\newtheorem{proposition}[subsection]{Proposition}
\newtheorem{lemma}[subsection]{Lemma}

\theoremstyle{definition}
\newtheorem{definition}[subsection]{Definition}
\newtheorem{example}[subsection]{Example}
\newtheorem{exercise}[subsection]{Exercise}
\newtheorem{situation}[subsection]{Situation}

\theoremstyle{remark}
\newtheorem{remark}[subsection]{Remark}
\newtheorem{remarks}[subsection]{Remarks}

\numberwithin{equation}{subsection}

% Macros
%
\def\lim{\mathop{\mathrm{lim}}\nolimits}
\def\colim{\mathop{\mathrm{colim}}\nolimits}
\def\Spec{\mathop{\mathrm{Spec}}}
\def\Hom{\mathop{\mathrm{Hom}}\nolimits}
\def\Ext{\mathop{\mathrm{Ext}}\nolimits}
\def\SheafHom{\mathop{\mathcal{H}\!\mathit{om}}\nolimits}
\def\SheafExt{\mathop{\mathcal{E}\!\mathit{xt}}\nolimits}
\def\Sch{\mathit{Sch}}
\def\Mor{\mathop{\mathrm{Mor}}\nolimits}
\def\Ob{\mathop{\mathrm{Ob}}\nolimits}
\def\Sh{\mathop{\mathit{Sh}}\nolimits}
\def\NL{\mathop{N\!L}\nolimits}
\def\CH{\mathop{\mathrm{CH}}\nolimits}
\def\proetale{{pro\text{-}\acute{e}tale}}
\def\etale{{\acute{e}tale}}
\def\QCoh{\mathit{QCoh}}
\def\Ker{\mathop{\mathrm{Ker}}}
\def\Im{\mathop{\mathrm{Im}}}
\def\Coker{\mathop{\mathrm{Coker}}}
\def\Coim{\mathop{\mathrm{Coim}}}

% Boxtimes
%
\DeclareMathSymbol{\boxtimes}{\mathbin}{AMSa}{"02}

%
% Macros for moduli stacks/spaces
%
\def\QCohstack{\mathcal{QC}\!\mathit{oh}}
\def\Cohstack{\mathcal{C}\!\mathit{oh}}
\def\Spacesstack{\mathcal{S}\!\mathit{paces}}
\def\Quotfunctor{\mathrm{Quot}}
\def\Hilbfunctor{\mathrm{Hilb}}
\def\Curvesstack{\mathcal{C}\!\mathit{urves}}
\def\Polarizedstack{\mathcal{P}\!\mathit{olarized}}
\def\Complexesstack{\mathcal{C}\!\mathit{omplexes}}
% \Pic is the operator that assigns to X its picard group, usage \Pic(X)
% \Picardstack_{X/B} denotes the Picard stack of X over B
% \Picardfunctor_{X/B} denotes the Picard functor of X over B
\def\Pic{\mathop{\mathrm{Pic}}\nolimits}
\def\Picardstack{\mathcal{P}\!\mathit{ic}}
\def\Picardfunctor{\mathrm{Pic}}
\def\Deformationcategory{\mathcal{D}\!\mathit{ef}}

% --- appended 2026-10-12 (Chief): the project own definitions of the reference/slogan/history
% display environments, copied mechanically from the shared Stacks preamble lines 8-12
% (\usepackage{verbatim} plus the three \newenvironment lines, which the chapter preamble keeps
% verbatim from the project preamble). Nothing added, removed or redefined: these environments
% carry no mathematical content and the project itself typesets them as comments. Without these
% three lines a counted statement whose own text uses them fails with Environment slogan undefined.
\usepackage{verbatim}
\newenvironment{reference}{\comment}{\endcomment}
\newenvironment{slogan}{\comment}{\endcomment}
\newenvironment{history}{\comment}{\endcomment}
\usepackage{cleveref}
\crefname{originaldefinition}{original definition}{original definitions}
\Crefname{originaldefinition}{Original definition}{Original definitions}
\makeatletter
\newlabel{constructions-definition-projective-bundle}{{01OB}{source}{}{sourcepointer-1}{}}
\newlabel{constructions-definition-projective-bundle@cref}{{[originaldefinition][0][]01OB}{[1][0][]source}{}{}{}}
\makeatother
\crefname{originallemma}{original lemma}{original lemmas}
\Crefname{originallemma}{Original lemma}{Original lemmas}
\makeatletter
\newlabel{constructions-lemma-apply-relative}{{01O4}{source}{}{sourcepointer-2}{}}
\newlabel{constructions-lemma-apply-relative@cref}{{[originallemma][0][]01O4}{[1][0][]source}{}{}{}}
\makeatother
\crefname{originallemma}{original lemma}{original lemmas}
\Crefname{originallemma}{Original lemma}{Original lemmas}
\makeatletter
\newlabel{constructions-lemma-relative-proj-separated}{{01O2}{source}{}{sourcepointer-3}{}}
\newlabel{constructions-lemma-relative-proj-separated@cref}{{[originallemma][0][]01O2}{[1][0][]source}{}{}{}}
\makeatother
\crefname{originallemma}{original lemma}{original lemmas}
\Crefname{originallemma}{Original lemma}{Original lemmas}
\makeatletter
\newlabel{schemes-lemma-section-immersion}{{01KT}{source}{}{sourcepointer-4}{}}
\newlabel{schemes-lemma-section-immersion@cref}{{[originallemma][0][]01KT}{[1][0][]source}{}{}{}}
\makeatother
\title{Sections of a projective bundle and surjections onto invertible modules}
\author{The Stacks Project Authors}
\date{stacks-project:divisors.tex:local-snapshot-20261012; GNU FDL 1.2 or later}
\begin{document}
\maketitle
\setcounter{subsection}{1}
\noindent\textit{Editorial note.} All delivered mathematical text is the author's own, copied from the pinned Stacks Project chapter file divisors.tex; no statement or proof step was rewritten, reordered or removed.\par\medskip
\noindent\textit{Editorial note.} Carrier window: the statement and the proof both lie inside the window 8000-10000 granted to this lane by CONTROL 8.83; the rest of the chapter is untouched.\par\medskip
\noindent\textit{Editorial note.} The statement contains two matrix displays; they are part of the counted statement span and are preserved verbatim.\par\medskip
\noindent\textit{Editorial note.} The references of the counted spans are recorded as non-load-bearing navigation pointers; the argument that consumes them stays inside the counted proof span.\par\medskip
\noindent\textit{Editorial note.} Preamble: sources/stacks-divisors/retained-preamble.tex (88 lines, verbatim present).\par\medskip
\noindent\textit{Allowed standard prerequisites (editorial).}\begin{itemize}
\item The projective bundle as relative Proj of the symmetric algebra. Constructions, Definition constructions-definition-projective-bundle, used at proof lines 8219-8221. Identifies P(E) with the relative Proj that is used throughout.
\item Description of the functor of points of the relative Proj. Constructions, Lemma constructions-lemma-apply-relative, used at proof lines 8221-8224. Produces the bijection between sections and surjections.
\item The relative Proj is separated and a section of a separated morphism is a closed immersion. Constructions, Lemma constructions-lemma-relative-proj-separated and Schemes, Lemma schemes-lemma-section-immersion, used at proof lines 8225-8227. Show that the section is a closed immersion.
\end{itemize}\medskip
\ifdefined\hypertarget\hypertarget{sourcepointer-1}{}\fi
\noindent\textit{Source reference.} Original definition 01OB: focused-seeds/AGP/sources/stacks-snapshot-20261011/tags.txt, source lines 2188--2188. Navigation pointer into the Stacks Project; the standard fact it supplies is applied inside the counted proof span and the argument resting on it stays there.\par\medskip
\ifdefined\hypertarget\hypertarget{sourcepointer-2}{}\fi
\noindent\textit{Source reference.} Original lemma 01O4: focused-seeds/AGP/sources/stacks-snapshot-20261011/tags.txt, source lines 2181--2181. Navigation pointer into the Stacks Project; the standard fact it supplies is applied inside the counted proof span and the argument resting on it stays there.\par\medskip
\ifdefined\hypertarget\hypertarget{sourcepointer-3}{}\fi
\noindent\textit{Source reference.} Original lemma 01O2: focused-seeds/AGP/sources/stacks-snapshot-20261011/tags.txt, source lines 2179--2179. Navigation pointer into the Stacks Project; the standard fact it supplies is applied inside the counted proof span and the argument resting on it stays there.\par\medskip
\ifdefined\hypertarget\hypertarget{sourcepointer-4}{}\fi
\noindent\textit{Source reference.} Original lemma 01KT: focused-seeds/AGP/sources/stacks-snapshot-20261011/tags.txt, source lines 2062--2062. Navigation pointer into the Stacks Project; the standard fact it supplies is applied inside the counted proof span and the argument resting on it stays there.\par\medskip
\begin{lemma}
\label{lemma-conormal-sheaf-section-projective-bundle}
Let $X$ be a scheme. Let $\mathcal{E}$ be a quasi-coherent
$\mathcal{O}_X$-module. There is a bijection
$$
\left\{
\begin{matrix}
\text{sections }\sigma\text{ of the } \\
\text{morphism } \mathbf{P}(\mathcal{E}) \to X
\end{matrix}
\right\}
\leftrightarrow
\left\{
\begin{matrix}
\text{surjections }\mathcal{E} \to \mathcal{L}\text{ where} \\
\mathcal{L}\text{ is an invertible }\mathcal{O}_X\text{-module}
\end{matrix}
\right\}
$$
In this case $\sigma$ is a closed immersion and there is a canonical
isomorphism
$$
\Ker(\mathcal{E} \to \mathcal{L})
\otimes_{\mathcal{O}_X} \mathcal{L}^{\otimes -1}
\longrightarrow
\mathcal{C}_{\sigma(X)/\mathbf{P}(\mathcal{E})}
$$
Both the bijection and isomorphism are compatible with base change.
\end{lemma}

\begin{proof}
Recall that $\pi : \mathbf{P}(\mathcal{E}) \to X$ is the relative proj of the
symmetric algebra on $\mathcal{E}$, see
Constructions, Definition \ref{constructions-definition-projective-bundle}.
Hence the descriptions of sections $\sigma$ follows immediately from
the description of the functor of points of $\mathbf{P}(\mathcal{E})$
in Constructions, Lemma \ref{constructions-lemma-apply-relative}.
Since $\pi$ is separated, any section is a closed immersion
(Constructions, Lemma \ref{constructions-lemma-relative-proj-separated} and
Schemes, Lemma \ref{schemes-lemma-section-immersion}).
Let $U \subset X$ be an affine open and $k \in \mathcal{E}(U)$ and
$s \in \mathcal{E}(U)$ be local sections such that $k$ maps to
zero in $\mathcal{L}$ and $s$ maps to a generator $\overline{s}$
of $\mathcal{L}$.
Then $f = k/s$ is a section of $\mathcal{O}_{\mathbf{P}(\mathcal{E})}$
defined in an open neighbourhood $D_+(s)$ of $s(U)$ in $\pi^{-1}(U)$.
Moreover, since $k$ maps to zero in $\mathcal{L}$ we see that
$f$ is a section of the ideal sheaf of $s(U)$ in $\pi^{-1}(U)$.
Thus we can take the image $\overline{f}$ of $f$ in
$\mathcal{C}_{\sigma(X)/\mathbf{P}(\mathcal{E})}(U)$.
We claim (1) that the image $\overline{f}$ depends only on the
sections $k$ and $\overline{s}$ and not on the choice of $s$
and (2) that we get an isomorphism over $U$ in this manner (see below).
However, once (1) and (2) are established, we see that
the construction is compatible with base change by $U' \to U$
where $U'$ is affine, which proves that these local maps glue
and are compatible with arbitrary base change.

\medskip\noindent
To prove (1) and (2) we make explicit what is going on.
Namely, say $U = \Spec(A)$ and say $\mathcal{E} \to \mathcal{L}$
corresponds to the map of $A$-modules $M \to N$. Then
$k \in K = \Ker(M \to N)$ and $s \in M$ maps to a generator $\overline{s}$
of $N$. Hence $M = K \oplus A s$. Thus
$$
\text{Sym}(M) = \text{Sym}(K)[s]
$$
Consider the identification $\text{Sym}(K) \to \text{Sym}(M)_{(s)}$
via the rule $g \mapsto g/s^n$ for $g \in \text{Sym}^n(K)$.
This gives an isomorphism $D_+(s) = \Spec(\text{Sym}(K))$ such
that $\sigma$ corresponds to the ring map $\text{Sym}(K) \to A$
mapping $K$ to zero. Via this isomorphism we see that the quasi-coherent
module corresponding to $K$ is identified with
$\mathcal{C}_{\sigma(U)/D_+(s)}$ proving (2).
Finally, suppose that $s' = k' + s$ for some $k' \in K$.
Then
$$
k/s' = (k/s) (s/s') = (k/s) (s'/s)^{-1} = (k/s) (1 + k'/s)^{-1}
$$
in an open neighbourhood of $\sigma(U)$ in $D_+(s)$. Thus we see that
$s'/s$ restricts to $1$ on $\sigma(U)$ and we see that $k/s'$ maps to
the same element of the conormal sheaf as does $k/s$ thereby proving (1).
\end{proof}

\end{document}
